% GENERATED FILE. Edit canonical lessons, then run npm run generate:books.
% canonical-source-hash: fnv1a64:e19202b5f7b6d7cf
% canonical-lessons: MR-C528-chand, MR-C528-virangula, MR-C528-phursat, MR-C528-karamnuk, MR-C528-dang-hone

\chapter{Hobbies and Leisure}
\label{ch:mr-hobbies-and-leisure}
\hlchaptermodality{528}

\begin{chapteropening}
\textbf{By the end of this chapter:} \emph{I can talk about hobbies, pastimes and free time.}

Complete the last lesson of chapter 528: I can talk about hobbies, pastimes and free time.
\end{chapteropening}

\section[\texorpdfstring{\mr{छंद} (chand)}{छंद (chand)}]{\mr{छंद} (chand) — a hobby}

\label{lesson:MR-C528-chand}



\begin{quote}
Before the new one: say the Marathi for to stir, then the Marathi for to get burnt.
\end{quote}

\subsection*{You'll want to know: \mr{छंद}}
\textbf{\mr{छंद}} — \emph{chand} — \textquotedblleft{}a hobby\textquotedblright{}.

\begin{quote}
\mr{माझ्या} \mr{बहिणीला} \mr{नाणी} \mr{गोळा} \mr{करण्याचा} \mr{छंद} \mr{आहे}.
\end{quote}

\begin{quote}
\emph{mājhyā bahiṇīlā nāṇī goḷā karṇyācā chand āhe.}
\end{quote}

\begin{quote}
My sister's hobby is collecting coins.
\end{quote}

Masculine.

The person takes -\mr{ला}: \mr{तिला} \mr{छंद} \mr{आहे}.

\begin{cousinweb}
\textbf{\mr{छंद}} (\emph{chand}) is a tatsama, taken from Sanskrit \emph{chanda}, \textquotedblleft{}pleasure, wish\textquotedblright{} (Monier-Williams).
\end{cousinweb}

\begin{grammarlens}[title={What you've built}]
One more word for describing what you lived and felt.
\end{grammarlens}

\subsection*{Your turn}
\begin{itemize}
\raggedright
  \item \emph{Say it:} \emph{chand}
  \item \emph{Say it:} \emph{chand}, once more
  \item \emph{Say it:} \emph{mājhyā bahiṇīlā nāṇī goḷā karṇyācā chand āhe.} — My sister's hobby is collecting coins.
  \item \emph{Recall:} say the Marathi for to stir, then the Marathi for to get burnt, then say \emph{chand} again
\end{itemize}

\begin{tcolorbox}[breakable,colback=teal!4,colframe=teal!35!black,title={Before you move on}]
What does \mr{छंद} mean? (\textquotedblleft{}A hobby\textquotedblright{}.) Say it once more.
\end{tcolorbox}
\section[\texorpdfstring{\mr{विरंगुळा} (viraṅguḷā)}{विरंगुळा (viraṅguḷā)}]{\mr{विरंगुळा} (viraṅguḷā) — a pastime; a pleasant break}

\label{lesson:MR-C528-virangula}



\begin{quote}
Before the new one: say the Marathi for to get burnt, then the Marathi for a hobby.
\end{quote}

\subsection*{You'll want to know: \mr{विरंगुळा}}
\textbf{\mr{विरंगुळा}} — \emph{viraṅguḷā} — \textquotedblleft{}a pastime; a pleasant break\textquotedblright{}.

\begin{quote}
\mr{संध्याकाळी} \mr{बागेत} \mr{काम} \mr{करणं} \mr{हा} \mr{माझा} \mr{विरंगुळा} \mr{आहे}.
\end{quote}

\begin{quote}
\emph{sandhyākāḷī bāget kām karṇaṁ hā mājhā viraṅguḷā āhe.}
\end{quote}

\begin{quote}
Working in the garden in the evening is my pastime.
\end{quote}

Masculine.

\begin{cousinweb}
\textbf{\mr{विरंगुळा}} (\emph{viraṅguḷā}) is a Marathi word for something that takes the mind off work; its origin is not settled.
\end{cousinweb}

\begin{grammarlens}[title={What you've built}]
One more word for describing what you lived and felt.
\end{grammarlens}

\subsection*{Your turn}
\begin{itemize}
\raggedright
  \item \emph{Say it:} \emph{viraṅguḷā}
  \item \emph{Say it:} \emph{viraṅguḷā}, once more
  \item \emph{Say it:} \emph{sandhyākāḷī bāget kām karṇaṁ hā mājhā viraṅguḷā āhe.} — Working in the garden in the evening is my pastime.
  \item \emph{Recall:} say the Marathi for to get burnt, then the Marathi for a hobby, then say \emph{viraṅguḷā} again
\end{itemize}

\begin{tcolorbox}[breakable,colback=teal!4,colframe=teal!35!black,title={Before you move on}]
What does \mr{विरंगुळा} mean? (\textquotedblleft{}A pastime; a pleasant break\textquotedblright{}.) Say it once more.
\end{tcolorbox}
\section[\texorpdfstring{\mr{फुरसत} (phursat)}{फुरसत (phursat)}]{\mr{फुरसत} (phursat) — free time; leisure}

\label{lesson:MR-C528-phursat}



\begin{quote}
Before the new one: say the Marathi for a hobby, then the Marathi for a pastime.
\end{quote}

\subsection*{You'll want to know: \mr{फुरसत}}
\textbf{\mr{फुरसत}} — \emph{phursat} — \textquotedblleft{}free time; leisure\textquotedblright{}.

\begin{quote}
\mr{आज} \mr{मला} \mr{अजिबात} \mr{फुरसत} \mr{नाही}.
\end{quote}

\begin{quote}
\emph{āj malā ajibāt phursat nāhī.}
\end{quote}

\begin{quote}
I have no free time at all today.
\end{quote}

Feminine.

\begin{cousinweb}
\textbf{\mr{फुरसत}} (\emph{phursat}) is a loan from Persian \emph{fursat}, \textquotedblleft{}leisure\textquotedblright{}, usually traced to Arabic \emph{furṣa}, \textquotedblleft{}occasion\textquotedblright{}.
\end{cousinweb}

\begin{grammarlens}[title={What you've built}]
One more word for describing what you lived and felt.
\end{grammarlens}

\subsection*{Your turn}
\begin{itemize}
\raggedright
  \item \emph{Say it:} \emph{phursat}
  \item \emph{Say it:} \emph{phursat}, once more
  \item \emph{Say it:} \emph{āj malā ajibāt phursat nāhī.} — I have no free time at all today.
  \item \emph{Recall:} say the Marathi for a hobby, then the Marathi for a pastime, then say \emph{phursat} again
\end{itemize}

\begin{tcolorbox}[breakable,colback=teal!4,colframe=teal!35!black,title={Before you move on}]
What does \mr{फुरसत} mean? (\textquotedblleft{}Free time; leisure\textquotedblright{}.) Say it once more.
\end{tcolorbox}
\section[\texorpdfstring{\mr{करमणूक} (karamṇūk)}{करमणूक (karamṇūk)}]{\mr{करमणूक} (karamṇūk) — entertainment}

\label{lesson:MR-C528-karamnuk}



\begin{quote}
Before the new one: say the Marathi for a pastime, then the Marathi for free time.
\end{quote}

\subsection*{You'll want to know: \mr{करमणूक}}
\textbf{\mr{करमणूक}} — \emph{karamṇūk} — \textquotedblleft{}entertainment\textquotedblright{}.

\begin{quote}
\mr{जत्रेत} \mr{मुलांच्या} \mr{करमणुकीसाठी} \mr{खूप} \mr{खेळ} \mr{होते}.
\end{quote}

\begin{quote}
\emph{jatret mulāñcyā karamṇukīsāṭhī khūp kheḷ hote.}
\end{quote}

\begin{quote}
At the fair there were lots of games to entertain the children.
\end{quote}

Feminine.

The oblique is \mr{करमणुकी}-: \mr{करमणुकीसाठी}.

\begin{cousinweb}
\textbf{\mr{करमणूक}} (\emph{karamṇūk}) is built on \mr{करमणे}, \textquotedblleft{}to be diverted, to pass the time\textquotedblright{}; its origin is not settled.
\end{cousinweb}

\begin{grammarlens}[title={What you've built}]
One more word for describing what you lived and felt.
\end{grammarlens}

\subsection*{Your turn}
\begin{itemize}
\raggedright
  \item \emph{Say it:} \emph{karamṇūk}
  \item \emph{Say it:} \emph{karamṇūk}, once more
  \item \emph{Say it:} \emph{jatret mulāñcyā karamṇukīsāṭhī khūp kheḷ hote.} — At the fair there were lots of games to entertain the children.
  \item \emph{Recall:} say the Marathi for a pastime, then the Marathi for free time, then say \emph{karamṇūk} again
\end{itemize}

\begin{tcolorbox}[breakable,colback=teal!4,colframe=teal!35!black,title={Before you move on}]
What does \mr{करमणूक} mean? (\textquotedblleft{}Entertainment\textquotedblright{}.) Say it once more.
\end{tcolorbox}
\section[\texorpdfstring{\mr{दंग} \mr{होणे} (daṅg hoṇe)}{दंग होणे (daṅg hoṇe)}]{\mr{दंग} \mr{होणे} (daṅg hoṇe) — to be absorbed}

\label{lesson:MR-C528-dang-hone}



\begin{quote}
Before the new one: say the Marathi for free time, then the Marathi for entertainment.
\end{quote}

\subsection*{You'll want to know: \mr{दंग} \mr{होणे}}
\textbf{\mr{दंग} \mr{होणे}} — \emph{daṅg hoṇe} — \textquotedblleft{}to be absorbed\textquotedblright{}.

\begin{quote}
\mr{पुस्तक} \mr{वाचण्यात} \mr{अरुण} \mr{इतका} \mr{दंग} \mr{झाला} \mr{की} \mr{जेवायला} \mr{विसरला}.
\end{quote}

\begin{quote}
\emph{pustak vācṇyāt aruṇ itkā daṅg jhālā kī jevāylā visarlā.}
\end{quote}

\begin{quote}
Arun was so absorbed in his book that he forgot to eat.
\end{quote}

Past (he …): \mr{तो} \mr{खेळात} \mr{दंग} \mr{झाला}, \emph{to kheḷāt daṅg jhālā}.

What absorbs takes -\mr{त}: \mr{खेळात} \mr{दंग} \mr{होणे}.

\begin{cousinweb}
\textbf{\mr{दंग} \mr{होणे}} (\emph{daṅg hoṇe}) is \mr{दंग}, usually traced to Persian \emph{dang}, \textquotedblleft{}astonished, lost\textquotedblright{}, with \mr{होणे}.
\end{cousinweb}

\begin{grammarlens}[title={What you've built}]
One more word for describing what you lived and felt.
\end{grammarlens}

\subsection*{Your turn}
\begin{itemize}
\raggedright
  \item \emph{Say it:} \emph{daṅg hoṇe}
  \item \emph{Say it:} \emph{daṅg hoṇe}, once more
  \item \emph{Say it:} \emph{pustak vācṇyāt aruṇ itkā daṅg jhālā kī jevāylā visarlā.} — Arun was so absorbed in his book that he forgot to eat.
  \item \emph{Recall:} say the Marathi for free time, then the Marathi for entertainment, then say \emph{daṅg hoṇe} again
\end{itemize}

\begin{tcolorbox}[breakable,colback=teal!4,colframe=teal!35!black,title={Before you move on}]
What does \mr{दंग} \mr{होणे} mean? (\textquotedblleft{}To be absorbed\textquotedblright{}.) Say it once more.
\end{tcolorbox}
